\section{Постановка задачи}

\subsection{Входные данные и обозначения}

Пусть установка содержит $N\geq2$ неподвижные камеры, а калибровочный объект наблюдается в $M$ положениях. Под одним положением понимается набор наблюдений, для которых объект и камеры сохраняют взаимное расположение. Для движущейся мишени требуется синхронизация; для неподвижной допустима последовательная съёмка при сохранении её позы.

Для камеры $i$ предполагаются известными матрица внутренних параметров $K_i$ и параметры дисторсии $\boldsymbol\delta_i$. В системе объекта $O$ заданы координаты характерных точек $\mathbf X_j$. Для видимых точек измерены координаты $\mathbf u_{ijs}$ на изображении камеры $i$ в положении объекта $s$. Не требуется, чтобы каждая камера видела все точки или все положения.

Примем систему первой камеры $C_0$ за мировую систему $W$. Обозначим через $T_i\in SE(3)$ преобразование из $W$ в систему камеры $C_i$, а через $P_s\in SE(3)$ --- преобразование из системы объекта в $W$. Тогда $T_0=I$, а модель наблюдения имеет вид
\begin{equation}
  \widehat{\mathbf u}_{ijs}
  =\pi_i\!\left(T_iP_s\widetilde{\mathbf X}_j\right),
  \label{eq:projection}
\end{equation}
где $\widetilde{\mathbf X}_j$ --- однородные координаты точки, а $\pi_i$ включает перспективную проекцию, $K_i$ и модель дисторсии. Для устранения неоднозначности масштаба геометрия объекта должна быть известна в метрических единицах.

\subsection{Искомый результат и условия определимости}

Необходимо оценить $T_i$ для всех камер и вспомогательные позы $P_s$. Преобразование точки из камеры $i$ в систему камеры $j$ после калибровки задаётся матрицей
\begin{equation}
  T_{j\leftarrow i}=T_jT_i^{-1}.
  \label{eq:relative}
\end{equation}
Вместе с матрицами должны сохраняться направление преобразования, единицы длины и сведения об использованных наблюдениях.

Связь камер с наблюдаемыми положениями объекта можно представить двудольным графом. Ребро соответствует наличию допустимых наблюдений элементов мишени. Граф должен быть связным; иначе взаимное положение отдельных компонент не определяется. Связность является необходимым, но не достаточным условием: также нужны информативные расположения точек и позы, позволяющие исключить вырожденные конфигурации и неоднозначность плоской мишени. Для выбранной ниже инициализации дополнительно требуется связная цепочка наблюдений, допускающих самостоятельную оценку позы через PnP.

Сами по себе разные видимые грани не препятствуют калибровке, если их геометрическая связь известна. Однако ошибки размеров объекта, углов между гранями и расположения меток могут систематически смещать оценки. Поэтому номинальная модель куба и фактическая точность его изготовления должны учитываться раздельно.

\subsection{Требования к экспериментальным данным}

До реализации необходимо уточнить число и модели камер; наличие изображений, карт глубины и показателей достоверности; единицы глубины и её физический смысл; внутренние параметры и преобразование между каналами; геометрию и разметку объекта; порядок съёмки и наличие эталонных измерений. В текущем описании эти характеристики не считаются установленными.

Базовая ветвь метода предполагает доступность изображений, на которых различимы опорные элементы объекта. Если имеются только карты глубины или облака точек, потребуется отдельная процедура выделения геометрии мишени, и применение описанного ниже PnP без дополнительных данных будет невозможно. Использование глубины в совместной оптимизации рассматривается как расширение базовой ветви.
